% ============================================
% Assistant Professor Cover Letter
% Florida International University
% ============================================

\documentclass[11pt,letterpaper]{letter}

\usepackage{geometry}
\usepackage{microtype}
\usepackage[hidelinks]{hyperref}

\geometry{left=25mm,right=25mm,top=25mm,bottom=25mm}

\signature{Decory Edwards}
\address{Department of Economics \\ Johns Hopkins University \\ Baltimore, MD 21218 \\ \texttt{dedwar65@jh.edu}}

\begin{document}

\begin{letter}{Search Committee \\
Department of Economics \\
Steven J. Green School of International and Public Affairs \\
Florida International University}

\opening{Dear Members of the Search Committee,}

I am writing to apply for the tenure-track Assistant Professor position in the Department of Economics at Florida International University. I am a Ph.D. candidate in Economics at Johns Hopkins University, advised by Professors Christopher D. Carroll, Nicholas W. Papageorge, and Stelios Fourakis, and I expect to complete my degree in 2026. My research and teaching interests lie in macroeconomics, inequality, and applied empirical methods with clear policy relevance.

My research agenda focuses on understanding wealth inequality through the lens of heterogeneous returns to wealth and household behavior. In my job market paper, I use microdata from the Survey of Consumer Finances to estimate persistent heterogeneity in portfolio returns and embed these estimates in a structural macroeconomic model. I show that allowing for return heterogeneity substantially improves the model’s ability to match observed wealth concentration and generates realistic implications for consumption behavior and policy analysis. This work connects micro-level financial behavior to aggregate outcomes in a way that is empirically grounded and directly relevant for macroeconomic policy.

In related work, I use data from the Health and Retirement Study to study how trust and beliefs shape household investment outcomes. I find evidence of a nonlinear relationship between trust and realized returns, highlighting a behavioral channel through which social attitudes influence economic inequality. Together, my research emphasizes careful measurement, transparent empirical methods, and questions that sit at the intersection of inequality, development, and public policy.

\textbf{Teaching.} My teaching experience centers on undergraduate macroeconomics and applied policy courses. I have led weekly discussion sections and held regular office hours for Principles of Macroeconomics, Intermediate Macroeconomic Theory, and an upper-level macroeconomic policy seminar. My teaching philosophy emphasizes clarity, structure, and applied intuition. I aim to help students connect theory to data and policy by working through real examples and empirical evidence. I am prepared to teach undergraduate and graduate courses in macroeconomics, applied microeconomics, and quantitative methods.

Florida International University’s location, mission, and institutional strengths make it a particularly strong fit for my work. The Steven J. Green School of International and Public Affairs emphasizes empirical research that engages directly with policy-relevant questions in development, health, the environment, and global inequality. My research aligns naturally with this mission, both in its use of large-scale microdata and in its focus on distributional outcomes that matter for economic mobility and social welfare. FIU’s international orientation and strong graduate programs also provide an ideal environment for collaboration and for mentoring students interested in applied research careers in policy institutions, international organizations, and academia. I would be enthusiastic about contributing to the department’s research profile while teaching students to critically evaluate data, empirical methods, and policy tradeoffs.

I would welcome the opportunity to contribute to FIU’s teaching, research, and service missions as part of a dynamic, policy-oriented economics department. Thank you for your time and consideration. I look forward to the possibility of discussing my application with you.

\closing{Sincerely,}

\end{letter}
\end{document}
